\chapter{Homotopy and Covering Spaces}\label{homotopy}
\chapterstatus{complete}

\section{Introduction}\label{homotopy:section-introduction}

Homotopy theory studies spaces up to continuous deformation. Its first invariant is the fundamental group, which we
compute for the circle using covering spaces (Theorem~\ref{homotopy:theorem-circle}) and for spaces glued from simpler
pieces using the Seifert--van Kampen theorem (Theorem~\ref{homotopy:theorem-van-kampen}). Covering spaces of a
reasonable space are classified by subgroups of its fundamental group (Theorem~\ref{homotopy:theorem-classification}),
in close analogy with Galois theory. We then introduce CW complexes and the higher homotopy groups. Our main references
are \cite{hatcher-at} and \cite{spanier-at}.

\noindent\textbf{Prerequisites.} Chapter~\ref{general-topology} (General Topology) and Chapter~\ref{groups} (Groups).

Throughout, $I = [0, 1]$, $D^n = \{x \in \RR^n : \|x\|_2 \leq 1\}$ is the closed unit ball and $S^{n-1} = \partial D^n$ the unit sphere.
A \emph{map} of spaces means a continuous map.

\section{Homotopy}\label{homotopy:section-homotopy}

\begin{definition}\label{homotopy:definition-homotopy}
A \emph{homotopy} between maps $f, g \colon X \to Y$ is a map $H \colon X \times I \to Y$ with $H(x, 0) = f(x)$ and $H(x, 1) = g(x)$; we write
$f \simeq g$ and $H_t = H(-, t)$. It is a homotopy \emph{relative to} $A \subseteq X$ if $H_t(a) = f(a)$ for all $a \in A$ and $t \in I$. A map
$f \colon X \to Y$ is a \emph{homotopy equivalence} if there is $g \colon Y \to X$ with $gf \simeq \id_X$ and $fg \simeq \id_Y$; then $X$ and $Y$ are
\emph{homotopy equivalent}, $X \simeq Y$. A space is \emph{contractible} if it is homotopy equivalent to a point. A subspace $A \subseteq X$ is a
\emph{retract} if there is a map $r \colon X \to A$ with $r|_A = \id_A$, and a \emph{deformation retract} if moreover $ir \simeq \id_X$ relative to $A$,
where $i \colon A \to X$ is the inclusion.
\end{definition}

\begin{lemma}\label{homotopy:lemma-homotopy-relation}
Homotopy (relative to $A$) is an equivalence relation on maps $X \to Y$, compatible with composition: if $f \simeq f'$ and $g \simeq g'$, then
$gf \simeq g'f'$. Hence spaces and homotopy classes of maps form a category $\mathrm{h}\cat{Top}$, and homotopy equivalences are its isomorphisms.
\end{lemma}

\begin{proof}
Reflexivity: $H(x, t) = f(x)$. Symmetry: $H(x, 1 - t)$. Transitivity: if $H$ goes from $f$ to $g$ and $K$ from $g$ to $h$, the map equal to $H(x, 2t)$ for
$t \leq 1/2$ and $K(x, 2t - 1)$ for $t \geq 1/2$ is continuous by the pasting lemma (Lemma~\ref{general-topology:lemma-pasting}). If $H$ goes from $f$ to
$f'$ and $K$ from $g$ to $g'$, then $(x, t) \mapsto K(H(x, t), t)$ goes from $gf$ to $g'f'$; it is continuous as the composite of
$(x, t) \mapsto (H(x, t), t)$ and $K$. Relative homotopies are preserved by these constructions.
\end{proof}

\begin{example}\label{homotopy:example-contractible}
\begin{enumerate}
\item A convex subset $C \subseteq \RR^n$ is contractible: for $c \in C$, the \emph{straight-line homotopy} $H(x, t) = (1 - t)x + tc$ goes from $\id_C$ to the
constant map, and it is continuous by Proposition~\ref{general-topology:proposition-product}(5). More generally, any two maps $f, g \colon X \to C$ are
homotopic via $(1 - t)f + tg$, relative to the set where they agree.
\item $S^{n-1}$ is a deformation retract of $\RR^n \setminus \{0\}$ and of $D^n \setminus \{0\}$, via $H(x, t) = (1 - t)x + tx/\|x\|$, which never vanishes.
\end{enumerate}
\end{example}

\begin{lemma}\label{homotopy:lemma-quotient-times-interval}
If $q \colon X \to Y$ is a quotient map and $K$ is a compact Hausdorff space, then $q \times \id_K \colon X \times K \to Y \times K$ is a quotient map. If
$U \subseteq Y$ is open, then $q^{-1}(U) \to U$ is a quotient map.
\end{lemma}

\begin{proof}
Let $W \subseteq Y \times K$ with $W' = (q \times \id)^{-1}(W)$ open, and let $(y_0, k_0) \in W$ with $y_0 = q(x_0)$. The set
$\{k : (x_0, k) \in W'\}$ is open and contains $k_0$; since $K$ is regular (Proposition~\ref{general-topology:proposition-normal-examples}), there is an open
$N \ni k_0$ whose closure $\overline{N}$ lies in it. Let $U' = \{x \in X : \{x\} \times \overline{N} \subseteq W'\}$, so $x_0 \in U'$. The set $U'$ is open: for
$x \in U'$, each $(x, k)$ with $k \in \overline{N}$ has a basic neighbourhood $M_k \times O_k \subseteq W'$, finitely many $O_k$ cover the compact set
$\overline{N}$, and the intersection $M$ of the corresponding $M_k$ satisfies $M \times \overline{N} \subseteq W'$. It is saturated: if $q(x) = q(x')$ and
$x \in U'$, then $(x', k)$ and $(x, k)$ have the same image, and $W'$ is a union of fibres of $q \times \id$. So $q(U')$ is open, and
$(y_0, k_0) \in q(U') \times N \subseteq W$. Hence $W$ is open. For the second statement, if $V \subseteq U$ and $q^{-1}(V)$ is open in $q^{-1}(U)$, then it is
open in $X$, so $V$ is open in $Y$.
\end{proof}

\section{The fundamental group}\label{homotopy:section-fundamental-group}

\begin{definition}\label{homotopy:definition-paths}
For paths $\gamma, \delta \colon I \to X$ with $\gamma(1) = \delta(0)$, the \emph{concatenation} $\gamma * \delta$ is $\gamma(2s)$ for $s \leq 1/2$ and
$\delta(2s - 1)$ for $s \geq 1/2$ (first $\gamma$, then $\delta$), and the \emph{reverse} of $\gamma$ is $\bar{\gamma}(s) = \gamma(1 - s)$. We write
$\gamma \simeq \delta$ for homotopy relative to $\{0, 1\}$ and $[\gamma]$ for the class. A \emph{loop} at $x_0$ is a path with $\gamma(0) = \gamma(1) = x_0$,
and $c_{x_0}$ denotes the constant loop.
\end{definition}

\begin{lemma}\label{homotopy:lemma-groupoid}
Concatenation is compatible with $\simeq$, and for composable paths:
\begin{enumerate}
\item $\gamma \circ \varphi \simeq \gamma$ for every map $\varphi \colon I \to I$ with $\varphi(0) = 0$ and $\varphi(1) = 1$;
\item $(\gamma * \delta) * \epsilon \simeq \gamma * (\delta * \epsilon)$, and $c_{\gamma(0)} * \gamma \simeq \gamma \simeq \gamma * c_{\gamma(1)}$;
\item $\gamma * \bar{\gamma} \simeq c_{\gamma(0)}$ and $\bar\gamma * \gamma \simeq c_{\gamma(1)}$.
\end{enumerate}
\end{lemma}

\begin{proof}
If $H$ and $K$ are homotopies $\gamma \simeq \gamma'$ and $\delta \simeq \delta'$, then $H(-, t) * K(-, t)$ is a homotopy $\gamma * \delta \simeq \gamma' * \delta'$.
(1) $H(s, t) = \gamma((1 - t)\varphi(s) + ts)$. (2) Each side is the other composed with a piecewise affine $\varphi$ as in (1): for instance
$(\gamma * \delta) * \epsilon = (\gamma * (\delta * \epsilon)) \circ \varphi$ with $\varphi$ affine on $[0, \frac14]$, $[\frac14, \frac12]$, $[\frac12, 1]$
mapping them to $[0, \frac12]$, $[\frac12, \frac34]$, $[\frac34, 1]$. (3) $H(s, t) = \gamma(2s(1 - t))$ for $s \leq 1/2$ and
$H(s, t) = \gamma((2 - 2s)(1 - t))$ for $s \geq 1/2$ is a homotopy from $\gamma * \bar\gamma$ to $c_{\gamma(0)}$; apply this to $\bar\gamma$ for the second
statement.
\end{proof}

\begin{definition}\label{homotopy:definition-fundamental-group}
The \emph{fundamental group} $\pi_1(X, x_0)$ is the set of classes of loops at $x_0$ with the multiplication $[\gamma][\delta] = [\gamma * \delta]$, which is a
group by Lemma~\ref{homotopy:lemma-groupoid}, with identity $[c_{x_0}]$ and inverse $[\gamma]^{-1} = [\bar\gamma]$. (The group is written multiplicatively;
note that $[\gamma][\delta]$ means ``first $\gamma$''.) A map $f \colon (X, x_0) \to (Y, y_0)$ of pointed spaces induces a homomorphism
$f_*[\gamma] = [f \circ \gamma]$, and $(gf)_* = g_*f_*$, $\id_* = \id$. A path-connected space $X$ is \emph{simply connected} if $\pi_1(X, x_0)$ is trivial.
\end{definition}

\begin{lemma}\label{homotopy:lemma-change-basepoint}
A path $h$ from $x_0$ to $x_1$ induces an isomorphism $\beta_h \colon \pi_1(X, x_0) \to \pi_1(X, x_1)$, $[\gamma] \mapsto [\bar h * \gamma * h]$. If $H$ is a homotopy
from $f$ to $g \colon X \to Y$ and $h(t) = H(x_0, t)$, then $g_* = \beta_h \circ f_* \colon \pi_1(X, x_0) \to \pi_1(Y, g(x_0))$. Consequently a homotopy
equivalence induces isomorphisms on fundamental groups, and contractible spaces are simply connected.
\end{lemma}

\begin{proof}
$\beta_h$ is a homomorphism by Lemma~\ref{homotopy:lemma-groupoid}, with inverse $\beta_{\bar h}$. For the second statement let $\gamma$ be a loop at $x_0$ and
$F(s, t) = H(\gamma(s), t)$ on the square $I^2$. In the convex set $I^2$, the path along the bottom and then the right side and the path along the left side and
then the top have the same endpoints, so they are homotopic relative to the endpoints (Example~\ref{homotopy:example-contractible}(1)). Composing with
$F$ gives $(f \circ \gamma) * h \simeq h * (g \circ \gamma)$, that is, $g_*[\gamma] = [\bar h * (f \circ \gamma) * h]$.

Let $f$ be a homotopy equivalence with inverse $g$. Since $gf \simeq \id_X$, the composite
$\pi_1(X, x_0) \xrightarrow{f_*} \pi_1(Y, f(x_0)) \xrightarrow{g_*} \pi_1(X, gf(x_0))$ is $\beta_h$, a bijection. Since $fg \simeq \id_Y$, the composite
$\pi_1(Y, f(x_0)) \xrightarrow{g_*} \pi_1(X, gf(x_0)) \xrightarrow{f_*} \pi_1(Y, fgf(x_0))$ is also a bijection. So this $g_*$ is injective and surjective, hence
bijective, and then $f_* = g_*^{-1}(g_*f_*)$ is bijective.
\end{proof}

\begin{proposition}\label{homotopy:proposition-product}
$\pi_1(X \times Y, (x_0, y_0)) \cong \pi_1(X, x_0) \times \pi_1(Y, y_0)$ via the projections.
\end{proposition}

\begin{proof}
Loops and homotopies in $X \times Y$ are pairs of loops and homotopies in $X$ and $Y$ (Proposition~\ref{general-topology:proposition-product}(2)).
\end{proof}

\section{Covering spaces and the circle}\label{homotopy:section-coverings}

\begin{definition}\label{homotopy:definition-covering}
A \emph{covering map} is a surjective map $p \colon E \to B$ such that every $b \in B$ has an open neighbourhood $U$ that is \emph{evenly covered}: $p^{-1}(U)$ is a
disjoint union of open sets (the \emph{sheets} over $U$), each mapped homeomorphically onto $U$ by $p$. A \emph{lift} of a map $f \colon Y \to B$ is a map
$\tilde{f} \colon Y \to E$ with $p\tilde{f} = f$.
\end{definition}

\begin{lemma}[Unique lifting]\label{homotopy:lemma-unique-lifting}
Let $p \colon E \to B$ be a covering map and $Y$ connected. Two lifts of $f \colon Y \to B$ that agree at one point are equal.
\end{lemma}

\begin{proof}
Let $\tilde{f}_1, \tilde{f}_2$ be lifts, $y \in Y$, $U$ an evenly covered neighbourhood of $f(y)$, and $\tilde{U}_1, \tilde{U}_2$ the sheets containing
$\tilde{f}_1(y), \tilde{f}_2(y)$. On the open neighbourhood $N = \tilde{f}_1^{-1}(\tilde{U}_1) \cap \tilde{f}_2^{-1}(\tilde{U}_2)$ of $y$: if
$\tilde{f}_1(y) = \tilde{f}_2(y)$, then $\tilde{U}_1 = \tilde{U}_2$ and, $p$ being injective on it, $\tilde{f}_1 = \tilde{f}_2$ on $N$; if
$\tilde{f}_1(y) \neq \tilde{f}_2(y)$, then $\tilde{U}_1 \neq \tilde{U}_2$ (as $p$ is injective on sheets), the sheets are disjoint and $\tilde{f}_1 \neq \tilde{f}_2$
on $N$. So the set where the lifts agree is open and closed, and nonempty, hence all of $Y$.
\end{proof}

\begin{proposition}[Homotopy lifting]\label{homotopy:proposition-homotopy-lifting}
Let $p \colon E \to B$ be a covering map, $F \colon Y \times I \to B$ a homotopy and $\tilde{F}_0 \colon Y \to E$ a lift of $F(-, 0)$. There is a unique lift
$\tilde{F} \colon Y \times I \to E$ of $F$ with $\tilde{F}(-, 0) = \tilde{F}_0$. In particular (for $Y$ a point) every path in $B$ has a unique lift with a
given initial point.
\end{proposition}

\begin{proof}
Fix $y_0 \in Y$. Each $(y_0, t)$ has a neighbourhood $N_t \times J_t$ mapped by $F$ into an evenly covered set. By compactness of $I$ there are an open
$N \ni y_0$ and a partition $0 = t_0 < t_1 < \cdots < t_m = 1$ with $F(N \times [t_i, t_{i+1}]) \subseteq U_i$, $U_i$ evenly covered (intersect finitely many
$N_t$, and take a partition subordinate to a Lebesgue number of the cover $(J_t)$, Lemma~\ref{general-topology:lemma-lebesgue-number}). Suppose a lift
$\tilde{F}$ has been constructed on $N \times [0, t_i]$ (for $i = 0$ this is $\tilde{F}_0$ on $N \times \{0\}$). Let $\tilde{U}_i$ be the sheet over $U_i$
containing $\tilde{F}(y_0, t_i)$; replacing $N$ by the open set $\{y \in N : \tilde{F}(y, t_i) \in \tilde{U}_i\}$, define
$\tilde{F} = (p|_{\tilde{U}_i})^{-1} \circ F$ on $N \times [t_i, t_{i+1}]$; this agrees on $N \times \{t_i\}$ with the previous definition, and the pasting lemma
gives a lift on $N \times [0, t_{i+1}]$. After $m$ steps we have a lift on $N_{y_0} \times I$ for an open $N_{y_0} \ni y_0$.

For each $y$, the restrictions to $\{y\} \times I$ of two lifts agreeing at $(y, 0)$ coincide, by Lemma~\ref{homotopy:lemma-unique-lifting}. Hence the lifts on
the sets $N_{y_0} \times I$ agree on overlaps, glue to a continuous lift on $Y \times I$ (pasting lemma for the open cover), and this lift is unique.
\end{proof}

\begin{corollary}\label{homotopy:corollary-lifting-consequences}
Let $p \colon E \to B$ be a covering map, $e_0 \in E$ and $b_0 = p(e_0)$.
\begin{enumerate}
\item If $\gamma \simeq \gamma'$ are paths starting at $b_0$, their lifts starting at $e_0$ have the same endpoint and are homotopic.
\item $p_* \colon \pi_1(E, e_0) \to \pi_1(B, b_0)$ is injective, and its image consists of the classes of loops whose lift starting at $e_0$ is a loop.
\item The \emph{monodromy} $e \cdot [\gamma] = \tilde\gamma_e(1)$, where $\tilde\gamma_e$ is the lift of $\gamma$ starting at $e$, is a right action of
$\pi_1(B, b_0)$ on the fibre $p^{-1}(b_0)$. If $E$ is path-connected, the action is transitive and the stabiliser of $e_0$ is $p_*\pi_1(E, e_0)$; so the
fibre is in bijection with the set of right cosets of $p_*\pi_1(E, e_0)$.
\end{enumerate}
\end{corollary}

\begin{proof}
(1) Lift a homotopy $F$ from $\gamma$ to $\gamma'$ with $\tilde{F}(-, 0) = \tilde\gamma$. The paths $t \mapsto \tilde{F}(0, t)$ and $t \mapsto \tilde{F}(1, t)$ lift
constant paths, so they are constant by uniqueness; hence $\tilde{F}(-, 1)$ is the lift of $\gamma'$ starting at $e_0$ and has the same endpoint as $\tilde\gamma$.
(2) If $p \circ \tilde\gamma \simeq c_{b_0}$ for a loop $\tilde\gamma$, lifting the homotopy gives $\tilde\gamma \simeq c_{e_0}$ by (1). If the lift of $\gamma$ is a
loop $\tilde\gamma$, then $[\gamma] = p_*[\tilde\gamma]$; conversely the lift of $p \circ \tilde\gamma$ is $\tilde\gamma$, and (1) handles homotopic loops.
(3) The action is well defined by (1), and the lift of $\gamma * \delta$ starting at $e$ is $\tilde\gamma_e * \tilde\delta_{e \cdot [\gamma]}$, so
$e \cdot ([\gamma][\delta]) = (e \cdot [\gamma]) \cdot [\delta]$. If $\eta$ is a path from $e_0$ to $e$, then $e_0 \cdot [p \circ \eta] = e$. The stabiliser
statement is (2).
\end{proof}

\begin{proposition}\label{homotopy:proposition-circle-covering}
Let $Q = \{x \in \RR^2 : \|x\|_\infty = 1\}$ and let $P \colon \RR \to Q$ be the map of period $4$ that on $[0, 4]$ runs affinely through the vertices
$(1, -1)$, $(1, 1)$, $(-1, 1)$, $(-1, -1)$, $(1, -1)$ at the times $0, 1, 2, 3, 4$. Let $r \colon Q \to S^1$, $x \mapsto x/\|x\|_2$, and $p = r \circ P$. Then $p$ is a
covering map, $p(t) = p(t')$ if and only if $t - t' \in 4\ZZ$, and $p(t + 4) = p(t)$.
\end{proposition}

\begin{proof}
$P$ is continuous by the pasting lemma (applied to the closed intervals $[k, k+1]$ inside each bounded open interval), and $r$ is a homeomorphism with
inverse $x \mapsto x/\|x\|_\infty$ (Example~\ref{general-topology:example-circle}). By construction $P$ is injective on $[a, a + 4)$ for every $a$ and has
period $4$, which gives the second statement. For $a \in \RR$, the restriction $P|_{[a, a+4]}$ is a closed map (Proposition~\ref{general-topology:proposition-compact-basic}(4)).
If $V \subseteq (a, a + 4)$ is open, then $P(V) = Q \setminus P([a, a+4] \setminus V)$, because $P$ is surjective, injective on $(a, a+4)$ and identifies only
$a$ and $a + 4$; so $P(V)$ is open, and $P$ maps $(a, a+4)$ homeomorphically onto $Q \setminus \{P(a)\}$. Given $q \in Q$, choose $a$ with $P(a) \neq q$; then
$U = Q \setminus \{P(a)\}$ is an open neighbourhood of $q$ with $P^{-1}(U) = \bigsqcup_{k \in \ZZ} (a + 4k, a + 4k + 4)$, each interval mapped homeomorphically onto $U$.
Composing with $r$ proves that $p$ is a covering map.
\end{proof}

\begin{theorem}\label{homotopy:theorem-circle}
Let $b_0 = p(0)$ and $\omega(s) = p(4s)$. Then $\pi_1(S^1, b_0) \cong \ZZ$, with $[\omega]$ corresponding to $1$. Hence $\pi_1(S^1, x) \cong \ZZ$ for every
$x \in S^1$.
\end{theorem}

\begin{proof}
The space $\RR$ is contractible, hence simply connected (Lemma~\ref{homotopy:lemma-change-basepoint}), and path-connected. By
Corollary~\ref{homotopy:corollary-lifting-consequences}(3), $[\gamma] \mapsto 0 \cdot [\gamma]$ is a bijection from $\pi_1(S^1, b_0)$ onto $p^{-1}(b_0) = 4\ZZ$. If
the lifts of $\gamma$ and $\delta$ starting at $0$ end at $4m$ and $4n$, then $t \mapsto \tilde\delta_0(t) + 4m$ lifts $\delta$ (as $p$ has period $4$) and starts at $4m$, so the
lift of $\gamma * \delta$ ends at $4(m + n)$. So $[\gamma] \mapsto \tilde\gamma_0(1)/4$ is an isomorphism onto $\ZZ$, and $\tilde\omega_0(s) = 4s$ ends at $4$. The last
statement follows from Lemma~\ref{homotopy:lemma-change-basepoint}.
\end{proof}

\begin{remark}\label{homotopy:remark-exponential}
With the exponential function of Chapter~\ref{real-analysis}, the usual covering is $\RR \to S^1$, $t \mapsto (\cos 2\pi t, \sin 2\pi t)$. We used a
piecewise affine model so as not to depend on analysis. A map $f \colon S^1 \to S^1$ induces multiplication by an integer $\deg f$ on
$\pi_1(S^1) \cong \ZZ$ (after identifying basepoints by a path, which does not matter since $\ZZ$ is abelian), its \emph{degree}.
\end{remark}

\begin{corollary}[Brouwer fixed point theorem in dimension $2$]\label{homotopy:corollary-brouwer}
There is no retraction $D^2 \to S^1$, and every map $f \colon D^2 \to D^2$ has a fixed point.
\end{corollary}

\begin{proof}
If $r \colon D^2 \to S^1$ were a retraction with inclusion $i$, then $r_*i_* = \id$ on $\pi_1(S^1) \cong \ZZ$ would factor through $\pi_1(D^2) = 0$. If $f$ had no fixed
point, let $u(x) = x - f(x) \neq 0$ and $r(x) = x + \lambda(x)u(x)$, where
\[
\lambda(x) = \frac{-\langle x, u \rangle + \sqrt{\langle x, u \rangle^2 + \|u\|^2(1 - \|x\|^2)}}{\|u\|^2}
\]
is the non-negative solution of $\|x + \lambda u\|^2 = 1$. The square root is continuous on $[0, \infty)$ since $|\sqrt{a} - \sqrt{b}| \leq \sqrt{|a - b|}$, so $r$
is continuous, and $r(x) \in S^1$. If $\|x\| = 1$, then $\langle x, u \rangle = 1 - \langle x, f(x) \rangle \geq 0$ by Cauchy--Schwarz, so $\lambda(x) = 0$ and
$r(x) = x$. This contradicts the first statement.
\end{proof}

\begin{example}\label{homotopy:example-coverings}
\begin{enumerate}
\item For $n \geq 1$ the map $S^1 \to S^1$, $p(t) \mapsto p(nt)$, is well defined and is a covering map with $n$ sheets; it corresponds to the subgroup
$n\ZZ \subseteq \ZZ$ in Theorem~\ref{homotopy:theorem-classification}.
\item The projection $q \colon S^n \to \RR\PP^n$ is a two-sheeted covering: for $x \in S^n$ the open hemisphere $W = \{y \in S^n : \langle x, y \rangle > 0\}$ is
mapped injectively by $q$, $q^{-1}(q(W)) = W \sqcup (-W)$ is open, so $q(W)$ is open and $q|_W$, $q|_{-W}$ are homeomorphisms onto it ($q$ is open, since
$q^{-1}(q(V)) = V \cup (-V)$). For $n \geq 2$, $S^n$ is simply connected (Corollary~\ref{homotopy:corollary-spheres}), so
$\pi_1(\RR\PP^n) \cong \ZZ/2\ZZ$ by Corollary~\ref{homotopy:corollary-lifting-consequences}(3).
\end{enumerate}
\end{example}

\section{The Seifert--van Kampen theorem}\label{homotopy:section-van-kampen}

\begin{theorem}[Seifert--van Kampen]\label{homotopy:theorem-van-kampen}
Let $X = U \cup V$ with $U$, $V$ open and $U$, $V$, $U \cap V$ path-connected, and $x_0 \in U \cap V$. The homomorphism
\[
\Phi \colon \pi_1(U, x_0) *_{\pi_1(U \cap V, x_0)} \pi_1(V, x_0) \to \pi_1(X, x_0)
\]
induced by the inclusions (Proposition~\ref{groups:proposition-amalgam-pushout}) is an isomorphism.
\end{theorem}

\begin{proof}
Write $G$ for the amalgamated product and $\iota_U, \iota_V$ for the maps from $\pi_1(U)$, $\pi_1(V)$ to $G$; by Definition~\ref{groups:definition-amalgam}, $G$ is
generated by their images. For loops $\gamma_1, \ldots, \gamma_k$ at $x_0$ let $\gamma_1 \cdots \gamma_k$ denote the loop running through $\gamma_i$ on
$[(i-1)/k, i/k]$; it is homotopic to the iterated concatenation by Lemma~\ref{homotopy:lemma-groupoid}(1).

\emph{Surjectivity.} Let $\gamma$ be a loop at $x_0$. By the Lebesgue number lemma there is a partition $0 = s_0 < \cdots < s_m = 1$ such that each
$\gamma([s_{i-1}, s_i])$ lies in a set $A_i \in \{U, V\}$; merging adjacent intervals with the same label, we may assume $A_i \neq A_{i+1}$, so
$\gamma(s_i) \in U \cap V$ for $0 < i < m$. Choose paths $g_i$ in $U \cap V$ from $x_0$ to $\gamma(s_i)$, with $g_0, g_m$ constant. Then
$\gamma \simeq \prod_i (g_{i-1} * \gamma|_{[s_{i-1}, s_i]} * \bar g_i)$ (with the restriction reparametrised to $I$), a product of loops in $U$ or $V$.

\emph{Injectivity.} Let $\gamma_1, \ldots, \gamma_k$ and $\gamma'_1, \ldots, \gamma'_l$ be loops at $x_0$, each in $U$ or $V$ (label them $A_i$ and $A'_j$), with
$F \colon I \times I \to X$ a homotopy from $\gamma_1 \cdots \gamma_k$ to $\gamma'_1 \cdots \gamma'_l$ relative to $\{0, 1\}$. We show
$\prod_i \iota_{A_i}[\gamma_i] = \prod_j \iota_{A'_j}[\gamma'_j]$ in $G$. This proves that $\ker \Phi$ is trivial: an element of the kernel can be written as
$\prod_i \iota_{A_i}[\gamma_i]$, and $\gamma_1 \cdots \gamma_k$ is then homotopic to the constant loop, which we take as $\gamma'_1$ with $l = 1$.

Choose $N$ divisible by $k$ and $l$ such that each small square $R$ of the grid of mesh $1/N$ is mapped by $F$ into $U$ or $V$ (Lebesgue number lemma), and fix
a label $A_R \in \{U, V\}$ with $F(R) \subseteq A_R$. For each grid vertex $v$ let $S_v$ be the intersection of the labels of the squares containing $v$ and, for
$v$ on the bottom (top) edge in the domain of $\gamma_i$ ($\gamma'_j$), also of $A_i$ ($A'_j$). Then $S_v$ is $U$, $V$ or $U \cap V$, so it is path-connected, and
it contains $F(v)$. Choose a path $g_v$ in $S_v$ from $x_0$ to $F(v)$, constant if $F(v) = x_0$ and $v$ lies on the boundary of the square. For an edge $e$
from $v$ to $w$ let $\ell_e = g_v * F|_e * \bar g_w$, a loop in the intersection of the labels of the (one or two) squares containing $e$. Put
$\phi(e) = \iota_{A_R}[\ell_e]$ for a square $R$ containing $e$; this does not depend on $R$, because if two squares with different labels contain $e$, then
$\ell_e$ lies in $U \cap V$ and the two images agree in $G$ by the amalgamation relations. For an edge path $P$ let $\phi(P)$ be the product of the $\phi(e)$ along $P$.

For a square $R$ with lower left vertex $a$, lower right $b$, upper left $c$ and upper right $d$, all four loops $\ell_e$ of its edges lie in $A_R$, and
$F|_R$ gives, in the convex set $R$, a homotopy relative to endpoints in $A_R$ between $F|_{[a,b]} * F|_{[b,d]}$ and $F|_{[a,c]} * F|_{[c,d]}$. Hence
$[\ell_{ab}][\ell_{bd}] = [g_a * F|_{[a,b]} * F|_{[b,d]} * \bar g_d] = [\ell_{ac}][\ell_{cd}]$ in $\pi_1(A_R)$. So replacing the two lower-right edges of a square
by its two upper-left edges does not change $\phi(P)$, and moving square by square the path $P_0$ along the bottom and right sides of $I^2$ to the path $P_1$ along
the left and top sides gives $\phi(P_0) = \phi(P_1)$. On the vertical sides $F$ is constant at $x_0$ and all $g_v$ are constant, so those edges contribute
$1$. The bottom edges in the domain of $\gamma_i$ have loops $\ell_e$ in $A_i$, and their product in $\pi_1(A_i)$ is $[\gamma_i]$ (the paths $g_v$ cancel, and
$g_v$ is constant at the ends of the domain); since each $\ell_e$ also lies in $A_R$, $\iota_{A_R}[\ell_e] = \iota_{A_i}[\ell_e]$. Hence
$\phi(P_0) = \prod_i \iota_{A_i}[\gamma_i]$, and likewise $\phi(P_1) = \prod_j \iota_{A'_j}[\gamma'_j]$.
\end{proof}

\begin{corollary}\label{homotopy:corollary-spheres}
For $n \geq 2$ the sphere $S^n$ is simply connected. If $X = U \cup V$ as in Theorem~\ref{homotopy:theorem-van-kampen} with $U \cap V$ simply connected, then
$\pi_1(X) \cong \pi_1(U) * \pi_1(V)$; if $V$ is simply connected, then $\pi_1(U) \to \pi_1(X)$ is surjective with kernel the normal subgroup generated by the
image of $\pi_1(U \cap V)$.
\end{corollary}

\begin{proof}
Let $U = S^n \setminus \{N\}$ and $V = S^n \setminus \{-N\}$. By stereographic projection (Example~\ref{general-topology:example-stereographic}), $U$ and $V$ are
homeomorphic to $\RR^n$, hence contractible, and $U \cap V$ is homeomorphic to $\RR^n \setminus \{0\}$, which is path-connected for $n \geq 2$. So
$\pi_1(S^n) \cong 1 *_{\pi_1(U \cap V)} 1 = 1$. The other statements follow from the presentation of the amalgamated product
(Definition~\ref{groups:definition-amalgam}): the relations from $\pi_1(U \cap V)$ are trivial in the first case, and in the second they identify the image of
each $h \in \pi_1(U \cap V)$ in $\pi_1(U)$ with $1$.
\end{proof}

\begin{example}\label{homotopy:example-wedge}
Let $X = S^1 \vee S^1$ be the wedge of two circles, the quotient of $S^1 \sqcup S^1$ identifying the two basepoints $b_0$ to a point $x_0$. Let $U$ ($V$) be
the union of the first (second) circle with the second (first) circle minus the point $-b_0$. The open arc $S^1 \setminus \{-b_0\}$ deformation retracts onto
$b_0$ (it is homeomorphic to an open interval by Proposition~\ref{homotopy:proposition-circle-covering}), so $U$ and $V$ deformation retract onto the
circles (Lemma~\ref{homotopy:lemma-quotient-times-interval} ensures continuity of the homotopies on the quotient), and $U \cap V$ is contractible. By
Corollary~\ref{homotopy:corollary-spheres}, $\pi_1(S^1 \vee S^1) \cong \ZZ * \ZZ$, the free group on two generators. It is not abelian, although
$\pi_1(S^1 \times S^1) \cong \ZZ^2$ by Proposition~\ref{homotopy:proposition-product}.
\end{example}

\begin{proposition}[Attaching cells]\label{homotopy:proposition-attaching-cells}
Let $X$ be path-connected, $n \geq 2$, $\varphi \colon S^{n-1} \to X$ a map, and $Y = X \cup_\varphi D^n$ the space obtained by attaching an $n$-cell
(Example~\ref{general-topology:example-quotients}). Let $z_0 \in S^{n-1}$ and $x_0 = \varphi(z_0)$. The inclusion induces a surjection
$\pi_1(X, x_0) \to \pi_1(Y, x_0)$. For $n = 2$ its kernel is the normal subgroup generated by $[\varphi \circ \omega]$, where $\omega$ is the loop of
Theorem~\ref{homotopy:theorem-circle} based at $z_0$; for $n \geq 3$ it is an isomorphism.
\end{proposition}

\begin{proof}
Let $q \colon X \sqcup D^n \to Y$ be the quotient map. For $C \subseteq X$ closed, $q^{-1}(q(C)) = C \sqcup \varphi^{-1}(C)$ is closed, so $q|_X$ is an embedding onto a
closed subspace, which we identify with $X$; similarly $q$ maps the open ball $\operatorname{int} D^n$ homeomorphically onto an open subset. Let
$V = q(\operatorname{int} D^n)$, which is contractible, and $U = Y \setminus \{q(0)\}$. On $q^{-1}(U) \times I = (X \times I) \sqcup ((D^n \setminus \{0\}) \times I)$ define
$H$ as the projection on $X \times I$ and as $q((1 - t)z + tz/\|z\|)$ on the second part; it respects the identifications, so by
Lemma~\ref{homotopy:lemma-quotient-times-interval} it induces a homotopy on $U \times I$, which is a deformation retraction of $U$ onto $X$. The set
$U \cap V$ is homeomorphic to $\operatorname{int} D^n \setminus \{0\}$, which is path-connected and deformation retracts onto the sphere of radius $1/2$; so
$\pi_1(U \cap V)$ is $\ZZ$ for $n = 2$ (generated by $c(s) = q(\omega(s)/2)$) and trivial for $n \geq 3$ (Corollary~\ref{homotopy:corollary-spheres}).

Work at the basepoint $y_1 = q(z_0/2) \in U \cap V$. By Corollary~\ref{homotopy:corollary-spheres}, $\pi_1(U, y_1) \to \pi_1(Y, y_1)$ is surjective with kernel the
normal subgroup generated by the image of $\pi_1(U \cap V, y_1)$. Transport to $x_0$ along the path $\delta(t) = H(y_1, t)$ from $y_1$ to $x_0$: the inclusion
$X \to U$ is a homotopy equivalence, and by Lemma~\ref{homotopy:lemma-change-basepoint} applied to the homotopy $H$, the class $[c]$ corresponds to
$[\bar\delta * c * \delta] = [H_1 \circ c] = [\varphi \circ \omega]$ in $\pi_1(X, x_0)$. Since the change-of-basepoint isomorphisms $\beta_\delta$ in $U$ and $Y$ are
compatible with the inclusion, the statement follows.
\end{proof}

\section{Classification of covering spaces}\label{homotopy:section-classification}

In this section $B$ is path-connected and locally path-connected. Then a covering space $E$ is locally path-connected, being locally homeomorphic to $B$,
and if $E$ is connected it is path-connected (Proposition~\ref{general-topology:proposition-components}).

\begin{proposition}[Lifting criterion]\label{homotopy:proposition-lifting-criterion}
Let $p \colon (E, e_0) \to (B, b_0)$ be a covering map, $Y$ path-connected and locally path-connected, and $f \colon (Y, y_0) \to (B, b_0)$ a map. A lift
$\tilde{f} \colon (Y, y_0) \to (E, e_0)$ exists if and only if $f_*\pi_1(Y, y_0) \subseteq p_*\pi_1(E, e_0)$, and it is then unique.
\end{proposition}

\begin{proof}
Necessity follows from $f_* = p_*\tilde{f}_*$, and uniqueness from Lemma~\ref{homotopy:lemma-unique-lifting}. Conversely, for $y \in Y$ choose a path $\eta$ from
$y_0$ to $y$ and let $\tilde{f}(y)$ be the endpoint of the lift of $f \circ \eta$ starting at $e_0$. If $\eta'$ is another such path, the loop
$h = (f \circ \eta') * (f \circ \bar\eta)$ represents an element of $f_*\pi_1(Y) \subseteq p_*\pi_1(E)$, so its lift starting at $e_0$ is a loop
(Corollary~\ref{homotopy:corollary-lifting-consequences}(2)). This lift is the lift of $f \circ \eta'$, ending at some $e'$, followed by the lift of $f \circ \bar\eta$
starting at $e'$, which therefore ends at $e_0$; reversing it gives the lift of $f \circ \eta$ starting at $e_0$, which ends at $e'$. So $\tilde{f}$ is well defined,
and $p\tilde{f} = f$.

Continuity: let $y \in Y$, $U$ an evenly covered neighbourhood of $f(y)$ and $\tilde U$ the sheet containing $\tilde f(y)$. Choose a path-connected open
$W \ni y$ with $f(W) \subseteq U$. For $y' \in W$, using $\eta * \sigma$ with $\sigma$ a path in $W$ from $y$ to $y'$, the lift of $f \circ \sigma$ starting at
$\tilde{f}(y)$ is $(p|_{\tilde U})^{-1} \circ f \circ \sigma$. So $\tilde{f} = (p|_{\tilde U})^{-1} \circ f$ on $W$.
\end{proof}

\begin{definition}\label{homotopy:definition-deck}
Two coverings $p \colon E \to B$ and $p' \colon E' \to B$ are \emph{isomorphic} if there is a homeomorphism $\tau \colon E \to E'$ with $p'\tau = p$. The
automorphisms of a covering form the group $\mathrm{Deck}(p)$ of \emph{deck transformations}. A covering with $E$ path-connected is \emph{normal} (or
Galois) if $\mathrm{Deck}(p)$ acts transitively on the fibres, and \emph{universal} if $E$ is simply connected.
\end{definition}

\begin{proposition}\label{homotopy:proposition-deck}
Let $p \colon E \to B$ be a covering with $E$ path-connected, $e_0 \in p^{-1}(b_0)$ and $H = p_*\pi_1(E, e_0)$.
\begin{enumerate}
\item For $e \in p^{-1}(b_0)$ there is a deck transformation $\tau$ with $\tau(e_0) = e$ if and only if $p_*\pi_1(E, e) = H$. For $e = e_0 \cdot [\gamma]$ one has
$p_*\pi_1(E, e) = [\gamma]^{-1}H[\gamma]$.
\item The covering is normal if and only if $H$ is a normal subgroup, and then $\mathrm{Deck}(p) \cong \pi_1(B, b_0)/H$. In particular, for a universal covering,
$\mathrm{Deck}(p) \cong \pi_1(B, b_0)$.
\end{enumerate}
\end{proposition}

\begin{proof}
(1) If $p_*\pi_1(E, e) = H$, the lifting criterion gives maps $\tau, \tau' \colon E \to E$ over $B$ with $\tau(e_0) = e$ and $\tau'(e) = e_0$; then $\tau'\tau$ and
$\tau\tau'$ are lifts of $p$ fixing a point, hence identities. The converse is clear. If $\tilde\gamma$ is the lift of $\gamma$ from $e_0$ to $e$, then
$[\tilde\alpha] \mapsto [\bar{\tilde\gamma} * \tilde\alpha * \tilde\gamma]$ is an isomorphism $\pi_1(E, e_0) \to \pi_1(E, e)$, whence the formula.

(2) The first statement follows from (1) and transitivity of monodromy (Corollary~\ref{homotopy:corollary-lifting-consequences}). A deck transformation commutes
with monodromy, $\tau(e \cdot [\gamma]) = \tau(e) \cdot [\gamma]$, since $\tau \circ \tilde\gamma_e$ lifts $\gamma$ from $\tau(e)$. For $H$ normal let $\Psi[\gamma]$ be
the deck transformation with $\Psi[\gamma](e_0) = e_0 \cdot [\gamma]$. Then
$\Psi([\gamma][\delta])(e_0) = (e_0 \cdot [\gamma]) \cdot [\delta] = \Psi[\gamma](e_0 \cdot [\delta]) = \Psi[\gamma]\Psi[\delta](e_0)$, so $\Psi$ is a
homomorphism (deck transformations agreeing at a point are equal). It is surjective by (1), and its kernel is the stabiliser $H$ of $e_0$.
\end{proof}

\begin{definition}\label{homotopy:definition-slsc}
$B$ is \emph{semilocally simply connected} if every point $b$ has a neighbourhood $U$ such that $\pi_1(U, b) \to \pi_1(B, b)$ is trivial.
\end{definition}

\begin{theorem}\label{homotopy:theorem-universal-cover}
If $B$ is path-connected, locally path-connected and semilocally simply connected, then $B$ has a universal covering.
\end{theorem}

\begin{proof}
Let $\tilde{B}$ be the set of classes $[\gamma]$ of paths starting at $b_0$ (up to homotopy relative to endpoints), and $p[\gamma] = \gamma(1)$. Let $\mathcal{U}$
be the set of path-connected open $U \subseteq B$ with $\pi_1(U) \to \pi_1(B)$ trivial (at any, hence every, basepoint in $U$). Every path-connected open subset of a
member of $\mathcal{U}$ is in $\mathcal{U}$, so by the hypotheses $\mathcal{U}$ is a basis of $B$. For $U \in \mathcal{U}$ and $\gamma(1) \in U$ let
$U_{[\gamma]} = \{[\gamma * \eta] : \eta \text{ a path in } U \text{ starting at } \gamma(1)\}$. Then $p \colon U_{[\gamma]} \to U$ is bijective: surjective as $U$ is
path-connected, and injective since for two paths $\eta, \eta'$ in $U$ with the same endpoints, $\eta * \bar\eta'$ is null-homotopic in $B$. If
$[\gamma'] \in U_{[\gamma]}$ then $U_{[\gamma']} = U_{[\gamma]}$. Hence if $[\gamma''] \in U_{[\gamma]} \cap V_{[\gamma']}$, choosing $W \in \mathcal{U}$ with
$\gamma''(1) \in W \subseteq U \cap V$ gives $[\gamma''] \in W_{[\gamma'']} \subseteq U_{[\gamma]} \cap V_{[\gamma']}$; so the sets $U_{[\gamma]}$ form a basis of a
topology on $\tilde{B}$ (Lemma~\ref{general-topology:lemma-basis}).

For $V \subseteq U$ in $\mathcal{U}$, $p$ maps the basic open sets $V_{[\gamma']} \subseteq U_{[\gamma]}$ onto $V$, and $p^{-1}(V) \cap U_{[\gamma]} = V_{[\gamma']}$ for
$[\gamma'] \in U_{[\gamma]}$ with $\gamma'(1) \in V$. So $p|_{U_{[\gamma]}}$ is a homeomorphism onto $U$. The set $p^{-1}(U)$ is the union of the sets $U_{[\gamma]}$ with
$\gamma(1) \in U$, and two of them are equal or disjoint. So $p$ is a covering map.

For $[\gamma] \in \tilde{B}$, let $\gamma_t(s) = \gamma(ts)$. The path $t \mapsto [\gamma_t]$ is continuous: near $t_0$, with $\gamma([t_0, t]) \subseteq U \in \mathcal{U}$,
it equals $(p|_{U_{[\gamma_{t_0}]}})^{-1} \circ \gamma$. It is the lift of $\gamma$ starting at $\tilde{b}_0 = [c_{b_0}]$ and ends at $[\gamma]$. So $\tilde{B}$ is
path-connected, and a loop $\gamma$ whose lift is a loop satisfies $[\gamma] = [c_{b_0}]$. By Corollary~\ref{homotopy:corollary-lifting-consequences}(2),
$p_*\pi_1(\tilde{B}) = 1$, and as $p_*$ is injective, $\tilde{B}$ is simply connected.
\end{proof}

\begin{theorem}[Classification of coverings]\label{homotopy:theorem-classification}
Let $B$ be path-connected, locally path-connected and semilocally simply connected, and $b_0 \in B$. Sending $p \colon (E, e_0) \to (B, b_0)$ to
$p_*\pi_1(E, e_0)$ is a bijection between isomorphism classes of pointed coverings with $E$ path-connected and subgroups of $\pi_1(B, b_0)$. Forgetting basepoints,
isomorphism classes of coverings with $E$ path-connected correspond to conjugacy classes of subgroups.
\end{theorem}

\begin{proof}
Injectivity: if $p_*\pi_1(E, e_0) = p'_*\pi_1(E', e'_0)$, the lifting criterion gives pointed maps $E \to E'$ and $E' \to E$ over $B$, whose composites are identities by
uniqueness of lifts.

Surjectivity: let $H \subseteq \pi_1(B, b_0)$ and $\tilde{B}$ be as in the proof of Theorem~\ref{homotopy:theorem-universal-cover}. Put $[\gamma] \sim [\gamma']$ if
$\gamma(1) = \gamma'(1)$ and $[\gamma * \bar\gamma'] \in H$, an equivalence relation, and let $E_H = \tilde{B}/{\sim}$. If $[\gamma] \sim [\gamma']$ then
$[\gamma * \eta] \sim [\gamma' * \eta]$, so the saturation of a basic set $U_{[\gamma]}$ is the union of the sets $U_{[\gamma']}$ with $[\gamma'] \sim [\gamma]$; hence the
projection $\tilde{B} \to E_H$ is open, and the images of the sets $U_{[\gamma]}$ are open, pairwise equal or disjoint, and mapped bijectively (as distinct points of
$U_{[\gamma]}$ have distinct endpoints), continuously and openly onto $U$. So $E_H \to B$ is a covering, and $E_H$ is path-connected. The lift of a loop $\gamma$ at $b_0$
starting at the class of $[c_{b_0}]$ is $t \mapsto$ class of $[\gamma_t]$, which is a loop if and only if $[\gamma] \in H$; by
Corollary~\ref{homotopy:corollary-lifting-consequences}(2) the subgroup of $E_H$ is $H$.

Changing the basepoint within the fibre replaces the subgroup by a conjugate, and every conjugate arises (Proposition~\ref{homotopy:proposition-deck}); an
isomorphism of unpointed coverings carries $e_0$ to some point of the fibre, which gives the last statement.
\end{proof}

\begin{counterexample}\label{homotopy:counterexample-hawaiian-earring}
Semilocal simple connectivity is necessary: a space with a simply connected covering $p \colon E \to B$ is semilocally simply connected, because a loop in an evenly
covered $U$ lifts to a loop in a sheet, which is null-homotopic in $E$. The \emph{Hawaiian earring}, the union of the circles of radius $1/n$ centred at $(1/n, 0)$ in
$\RR^2$, is path-connected and locally path-connected but not semilocally simply connected (every neighbourhood of the origin contains a whole circle, which is a
retract of the earring and hence not null-homotopic in it), so it has no universal covering \cite{hatcher-at}.
\end{counterexample}

\section{CW complexes}\label{homotopy:section-cw}

\begin{definition}\label{homotopy:definition-cw}
A \emph{CW complex} is a space $X$ with an increasing sequence of subspaces $X^0 \subseteq X^1 \subseteq \cdots$, its \emph{skeleta}, such that $X^0$ is discrete,
$X^n$ is obtained from $X^{n-1}$ by attaching a family of $n$-cells, $X^n = (X^{n-1} \sqcup \coprod_\alpha D^n_\alpha)/(z \sim \varphi_\alpha(z))$ for maps
$\varphi_\alpha \colon S^{n-1} \to X^{n-1}$, and $X = \bigcup_n X^n$ carries the final topology for the inclusions of the $X^n$ (a set is open if and only if its
intersection with every $X^n$ is open). The images $e^n_\alpha$ of the open balls are the \emph{open cells}, the maps $\Phi_\alpha \colon D^n_\alpha \to X$ are the
\emph{characteristic maps}, and $X$ is \emph{finite} if there are finitely many cells and \emph{finite-dimensional} if $X = X^n$ for some $n$. A \emph{subcomplex}
is a closed union of cells containing, with each cell, the image of its characteristic map; $(X, A)$ is then a \emph{CW pair}.
\end{definition}

\begin{example}\label{homotopy:example-cw}
\begin{enumerate}
\item $S^n$ has a CW structure with one $0$-cell and one $n$-cell attached by the constant map, since $D^n/S^{n-1} \cong S^n$. It also has a structure with two
cells in each dimension $0, \ldots, n$ (the hemispheres), compatible with the antipodal map.
\item Passing to the quotient by the antipodal map in the second structure gives $\RR\PP^n$ one cell in each dimension $0, \ldots, n$, with
$\RR\PP^n = \RR\PP^{n-1} \cup e^n$ attached along the double covering $S^{n-1} \to \RR\PP^{n-1}$. Similarly $\CC\PP^n = \CC\PP^{n-1} \cup e^{2n}$ has one cell in each
even dimension (Chapter~\ref{complex-manifolds}).
\item A graph is a $1$-dimensional CW complex. The torus is obtained from $S^1 \vee S^1$ by attaching a $2$-cell along the loop $aba^{-1}b^{-1}$, and the closed
orientable surface of genus $g$ from a wedge of $2g$ circles along $a_1b_1a_1^{-1}b_1^{-1} \cdots a_gb_ga_g^{-1}b_g^{-1}$. By
Proposition~\ref{homotopy:proposition-attaching-cells} and Example~\ref{homotopy:example-wedge} (extended to $2g$ circles), its fundamental group is
$\langle a_1, b_1, \ldots, a_g, b_g \mid \prod_i a_ib_ia_i^{-1}b_i^{-1} \rangle$.
\end{enumerate}
\end{example}

\begin{theorem}\label{homotopy:theorem-cw-properties}
Let $X$ be a CW complex.
\begin{enumerate}
\item $X$ is Hausdorff, and in fact paracompact, and it is locally contractible.
\item Every compact subset of $X$ lies in a finite subcomplex.
\item (Homotopy extension property) If $(X, A)$ is a CW pair, then $X \times \{0\} \cup A \times I$ is a deformation retract of $X \times I$; so every map
$X \times \{0\} \cup A \times I \to Y$ extends to $X \times I$.
\item The inclusion $X^2 \to X$ induces an isomorphism on fundamental groups (at a $0$-cell).
\end{enumerate}
\end{theorem}

These are \cite[Appendix, and Propositions 0.16 and 1.26]{hatcher-at}. Part (4) follows from Proposition~\ref{homotopy:proposition-attaching-cells}
applied to the finite subcomplexes containing a given loop or homotopy, using (2).

\begin{theorem}[Cellular approximation]\label{homotopy:theorem-cellular-approximation}
Every map $f \colon X \to Y$ of CW complexes is homotopic to a \emph{cellular} map ($f(X^n) \subseteq Y^n$ for all $n$), by a homotopy relative to any subcomplex on which
$f$ is already cellular.
\end{theorem}

This is \cite[Theorem 4.8]{hatcher-at}.

\section{Higher homotopy groups and fibrations}\label{homotopy:section-higher}

\begin{definition}\label{homotopy:definition-higher-homotopy}
For $n \geq 1$, $\pi_n(X, x_0)$ is the set of homotopy classes, relative to $\partial I^n$, of maps $f \colon (I^n, \partial I^n) \to (X, x_0)$, with
$[f] + [g] = [f *_1 g]$, where $(f *_1 g)(s) = f(2s_1, s_2, \ldots, s_n)$ for $s_1 \leq 1/2$ and $g(2s_1 - 1, s_2, \ldots, s_n)$ for $s_1 \geq 1/2$. We write the
operation additively for $n \geq 2$ and multiplicatively for $n = 1$, where it is the fundamental group. Also $\pi_0(X, x_0)$ is the pointed set of path components.
\end{definition}

\begin{proposition}\label{homotopy:proposition-higher-abelian}
$\pi_n(X, x_0)$ is a group, abelian for $n \geq 2$. A pointed map induces homomorphisms $f_*$, pointed homotopic maps induce the same homomorphisms, and
$\pi_n(X \times Y) \cong \pi_n(X) \times \pi_n(Y)$.
\end{proposition}

\begin{proof}
The proofs of Lemma~\ref{homotopy:lemma-groupoid} apply verbatim in the first coordinate, the other coordinates being carried along. For $n \geq 2$ define $*_2$ in the
same way using the second coordinate; it also induces a well-defined operation on $\pi_n$ with the same identity $e$, the class of the constant map. For maps,
$(a *_1 b) *_2 (c *_1 d) = (a *_2 c) *_1 (b *_2 d)$ holds exactly (both sides are $a, b, c, d$ on the four quarters of $I^2 \times I^{n-2}$, rescaled). Hence on
classes
\[
a *_2 b = (a *_1 e) *_2 (e *_1 b) = (a *_2 e) *_1 (e *_2 b) = a *_1 b, \quad
a *_1 b = (e *_1 a) *_2 (b *_1 e) = (e *_2 b) *_1 (a *_2 e) = b *_1 a
\]
(the Eckmann--Hilton argument). The remaining statements are clear.
\end{proof}

\begin{proposition}\label{homotopy:proposition-covering-higher}
A covering map $p \colon E \to B$ induces isomorphisms $\pi_n(E, e_0) \to \pi_n(B, b_0)$ for $n \geq 2$. In particular $\pi_n(S^1) = 0$ and
$\pi_n(\RR\PP^m) \cong \pi_n(S^m)$ for $n \geq 2$.
\end{proposition}

\begin{proof}
For $n \geq 2$, $\partial I^n$ is homeomorphic to $S^{n-1}$ (by $x \mapsto x/\|x\|_2$ after centring, as in Example~\ref{general-topology:example-circle}), hence
path-connected. Let $f \colon (I^n, \partial I^n) \to (B, b_0)$. The cube is simply connected and locally path-connected, so $f$ has a lift $\tilde{f}$ with
$\tilde{f}(0) = e_0$ (Proposition~\ref{homotopy:proposition-lifting-criterion}); $\tilde{f}(\partial I^n)$ is a connected subset of the fibre $p^{-1}(b_0)$, which
is discrete, so $\tilde{f}(\partial I^n) = e_0$. So $p_*$ is surjective. If $p \circ \tilde{f}$ is null-homotopic relative to $\partial I^n$, lift the homotopy
(Proposition~\ref{homotopy:proposition-homotopy-lifting}); the image of the connected set $\partial I^n \times I$ lies in the fibre, hence is $e_0$, and at time $1$ the
lift of the constant map is constant. So $p_*$ is injective. Apply this to $\RR \to S^1$ ($\RR$ is contractible) and $S^m \to \RR\PP^m$.
\end{proof}

\begin{definition}\label{homotopy:definition-fibration}
A map $p \colon E \to B$ has the \emph{homotopy lifting property} with respect to $Y$ if for every homotopy $F \colon Y \times I \to B$ and lift $\tilde{F}_0$ of $F_0$ there
is a lift $\tilde{F}$ of $F$ with $\tilde{F}(-, 0) = \tilde{F}_0$ (not necessarily unique). It is a \emph{Hurewicz fibration} if this holds for all $Y$, and a
\emph{Serre fibration} if it holds for all cubes $Y = I^k$. The fibre over $b_0$ is $F = p^{-1}(b_0)$. For a pair $(X, A)$ with $x_0 \in A$ and $n \geq 1$, the
\emph{relative homotopy group} $\pi_n(X, A, x_0)$ is the set of homotopy classes of maps $(I^n, \partial I^n, J^{n-1}) \to (X, A, x_0)$, where
$J^{n-1} = \overline{\partial I^n \setminus I^{n-1} \times \{0\}}$; it is a group for $n \geq 2$.
\end{definition}

\begin{theorem}[Long exact sequences]\label{homotopy:theorem-long-exact-sequences}
\begin{enumerate}
\item For a pair $(X, A)$ there is a long exact sequence
$\cdots \to \pi_n(A) \to \pi_n(X) \to \pi_n(X, A) \xrightarrow{\partial} \pi_{n-1}(A) \to \cdots \to \pi_0(X)$, with $\partial[f] = [f|_{I^{n-1} \times \{0\}}]$.
\item For a Serre fibration $p \colon E \to B$ with $B$ path-connected, fibre $F$ and $x_0 \in F$, the map $p_* \colon \pi_n(E, F, x_0) \to \pi_n(B, b_0)$ is an
isomorphism for $n \geq 1$, which gives a long exact sequence
\[
\cdots \to \pi_n(F) \to \pi_n(E) \to \pi_n(B) \to \pi_{n-1}(F) \to \cdots \to \pi_0(E) \to 0.
\]
\end{enumerate}
\end{theorem}

\begin{proof}[Sketch of proof]
(1) Exactness at each term is a direct manipulation of maps of cubes, based on the \emph{compression criterion}: a map $f \colon (I^n, \partial I^n, J^{n-1}) \to (X, A, x_0)$
represents $0$ if and only if it is homotopic relative to $\partial I^n$ to a map into $A$. (2) Surjectivity of $p_*$ lifts a map $(I^n, \partial I^n) \to (B, b_0)$,
viewed as a homotopy of maps $I^{n-1} \to B$, starting from the constant lift; injectivity lifts a homotopy in the same way, using the homotopy lifting property for
$I^{n-1} \times I$ relative to $J^{n-1}$ (which holds for Serre fibrations because $(I^n, J^{n-1}) \cong (I^n, I^{n-1} \times \{0\})$). Substituting into (1) for the pair
$(E, F)$ gives the sequence. The omitted verifications are in \cite[Theorems 4.3 and 4.41]{hatcher-at}.
\end{proof}

\begin{theorem}[Whitehead]\label{homotopy:theorem-whitehead}
A map between connected CW complexes that induces isomorphisms on all homotopy groups $\pi_n$, $n \geq 1$, is a homotopy equivalence.
\end{theorem}

This is \cite[Theorem 4.5]{hatcher-at}.

\begin{remark}\label{homotopy:remark-higher-computations}
Higher homotopy groups are hard to compute. By cellular approximation, $\pi_k(S^n) = 0$ for $k < n$ (a map $S^k \to S^n$ is homotopic to a cellular map, which misses the
open $n$-cell and so factors through a point). The \emph{Hopf fibration} $S^3 \to S^2 = \CC\PP^1$, $(z_1, z_2) \mapsto [z_1 : z_2]$, has fibre $S^1$, and its long exact
sequence gives $\pi_3(S^2) \cong \pi_3(S^3) \cong \ZZ$, although $H_3(S^2) = 0$. Homology (Chapter~\ref{singular-homology}) is much more computable, and the
Hurewicz theorem relates the two.
\end{remark}

\section{Exercises}\label{homotopy:section-exercises}

\begin{exercise}\label{homotopy:exercise-retract-injective}
Show that if $A \subseteq X$ is a retract, then $i_* \colon \pi_1(A, a) \to \pi_1(X, a)$ is injective, and if $A$ is a deformation retract it is an isomorphism.
\end{exercise}

\begin{solution}
$r_*i_* = (ri)_* = \id$, so $i_*$ is injective. For a deformation retract, $ir \simeq \id_X$ relative to $A$, so $i$ is a homotopy equivalence and
Lemma~\ref{homotopy:lemma-change-basepoint} applies (the path $h$ is constant since $a \in A$ is fixed).
\end{solution}

\begin{exercise}\label{homotopy:exercise-topological-group}
Let $G$ be a topological group (a group with continuous multiplication and inversion) with identity $1$. Show that $\pi_1(G, 1)$ is abelian, and that the group
law of $\pi_1$ coincides with $[\gamma][\delta] = [s \mapsto \gamma(s)\delta(s)]$.
\end{exercise}

\begin{solution}
Let $\gamma \cdot \delta$ denote the pointwise product. For loops, $(\gamma * \delta) \cdot (\gamma' * \delta') = (\gamma \cdot \gamma') * (\delta \cdot \delta')$ exactly, and
the constant loop is a two-sided unit for both operations on classes. The Eckmann--Hilton argument of Proposition~\ref{homotopy:proposition-higher-abelian} shows
that the two operations coincide and are commutative.
\end{solution}

\begin{exercise}\label{homotopy:exercise-sphere-to-circle}
Show that every map $S^n \to S^1$ with $n \geq 2$ is null-homotopic, and that every map $S^1 \to S^1$ of degree $0$ is null-homotopic.
\end{exercise}

\begin{solution}
In both cases $f_*\pi_1 = 0$ ($S^n$ is simply connected for $n \geq 2$; for degree $0$ by definition), so $f$ lifts to $\tilde{f} \colon S^n \to \RR$ by the lifting
criterion. Then $f = p \circ \tilde{f} \simeq p \circ c$ via $p \circ ((1 - t)\tilde f + tc)$, where $c$ is a constant.
\end{solution}

\begin{exercise}\label{homotopy:exercise-deck-circle}
Show that the deck transformation group of $p \colon \RR \to S^1$ consists of the translations $t \mapsto t + 4k$, $k \in \ZZ$, and deduce again
$\pi_1(S^1) \cong \ZZ$ from Proposition~\ref{homotopy:proposition-deck}.
\end{exercise}

\begin{exercise}\label{homotopy:exercise-borsuk-ulam}
(Borsuk--Ulam in dimension $2$.) Show that for every map $f \colon S^2 \to \RR^2$ there is $x$ with $f(x) = f(-x)$. (Hint: otherwise
$g(x) = (f(x) - f(-x))/\|f(x) - f(-x)\|$ is an odd map $S^2 \to S^1$; restrict to the equator and compare lifts to $\RR$ along a half circle.)
\end{exercise}

\begin{exercise}\label{homotopy:exercise-free-group-covering}
Find all connected two-sheeted coverings of $S^1 \vee S^1$, and describe the covering corresponding to the kernel of $\ZZ * \ZZ \to \ZZ/2\ZZ \times \ZZ/2\ZZ$.
\end{exercise}
